
% Definición de absorción
Absorción es la capacidad de un material para absorber o detener libres \textbf{Hidronio (\ce{H3O+})} y \textbf{peróxido de hidrógeno (\ce{H2O2})}, los cuales son biológicamente muy dañinos y pueden producir la muerte de la célula o puede mantenerla en estado latente donde se puede llegar a realizar la mitosis.

El núcleo es la parte más sensible a la radiación, a este nivel la radiación gamma lesionará los enlaces de la molécula de DNA, alterando el genoma con lo cual se encuentra la molécula de DNA, este genoma alterado puede llegar a producir un efecto en los cromosomas, que al dividirse producirá células alteradas con funciones diferentes a las de la célula madre. Esto se conoce como efecto genético de la radiación ionizante.
  
En general, podemos decir que la alteración de la radiación ionizante con las células, puede producir en ellas:

\begin{itemize}
    \item Aumento o disminución de volumen.
    \item Su muerte.
    \item Un estado latente.
    \item Mutaciones en las células hijas.
\end{itemize}

La \textbf{Radiosensibilidad} es la respuesta que las células ofrecen a la acción de la radiación, la cual es directamente proporcional a la capacidad reproductiva de las células y varía inversamente con el grado de diferenciación de los mismos, así como de el estado metabólico, su estado de división y su estado de nutrición.

Por lo tanto, las células que más rápidamente se dividen son las más radiosensibles (hematopoyéticas, piel, cristalino, epitelio, gastrointestinales, gonadas, etc.).

En función a los dos tipos de células que se encuentran en los organismos, los efectos biológicos de la radiación se clasifican en dos tipos: somáticos y hereditarios.

\textbf{Los somáticos} son aquellos que se manifiestan en el individuo que se ha expuesto a la radiación y cuyo número no indica que se deba al daño recibido en sus células somáticas o diploides; además este daño queda limitado solamente al individuo. Algunos ejemplos de este tipo de alteraciones son la esterilidad, disminución de células en la médula osea e inducción al cáncer.

\textbf{Los hereditarios o genéticos} son el resultado del daño recibido en las células gaméticas y sus efectos se presentarán en la descendencia de las células irradiadas. Un ejemplo serían los procesos mutagénicos en el \textbf{ADN}, los cuales serían heredados a sus descendientes.

Actualmente los efectos biológicos de la radiación, de acuerdo a su probabilidad de incidencia, se clasifican en dos tipos: efectos determinísticos (no estocásticos) y efectos estocásticos.

\textbf{Los efectos determinísticos} se definen como aquellos que se producen a partir de una dosis de umbral y aumentan en severidad dependiendo de la dosis recibida. El efecto se agudiza para dosis altas recibidas en un tiempo corto, estos efectos se derivan de la muerte celular.

La dosis umbral para algunos efectos determinísticos se presenta a continuación, cuando la irradiación es a cuerpo total.
